\documentclass[a4paper,12pt]{article}
\usepackage[utf8]{inputenc}
\usepackage[russian]{babel}
\usepackage[T2A]{fontenc}
\usepackage{amsmath, amsfonts, amssymb}
\usepackage{geometry}
\usepackage{graphicx}
\usepackage{tabularx}
\usepackage{booktabs}
\usepackage{multirow}
\usepackage{caption}

\geometry{left=2cm, right=2cm, top=2cm, bottom=2cm}

\begin{document}

% --- Титульный лист / Шапка ---
\begin{flushleft}
    Университет ИТМО \\
    Физико-технический мегафакультет \\
    Физический факультет \\
\end{flushleft}

\vspace{0.5cm}

\noindent
\begin{tabularx}{\textwidth}{|X|X|X|}
\hline
Группа: \textbf{P3210} & К работе допущен: & Студент: \textbf{Чжун Цзяцзюнь, Давидюк Антон} \\ \hline
Работа выполнена: & Преподаватель: \textbf{Сорокина Елена Константиновна} & Отчет принят: \\ \hline
\end{tabularx}

\vspace{1cm}

\begin{center}
    \Large \textbf{Рабочий протокол и отчет по лабораторной работе № 3.08 (2.12)} \\
    \large \textbf{ЭФФЕКТ ХОЛЛА В ПРИМЕСНЫХ ПОЛУПРОВОДНИКАХ}
\end{center}

\section{Цель работы}
Изучить эффект Холла в примесных полупроводниках. Ознакомиться с методом измерения концентрации и подвижности основных носителей тока в примесных полупроводниках с помощью эффекта Холла.

\section{Рабочие формулы}
\begin{enumerate}
    \item Постоянная Холла: $R_x = \frac{U_x \cdot d}{I \cdot B}$
    \item Концентрация носителей заряда: $n = \frac{1}{R_x \cdot e}$
    \item Удельная электропроводность: $\sigma = \frac{I \cdot L}{U \cdot a \cdot d}$
    \item Подвижность носителей: $\mu = \sigma \cdot R_x$
    \item Напряжение Холла: $U_x = \frac{U'_{34} - U''_{34}}{2}$ (для исключения побочных эффектов)
\end{enumerate}

\section{Результаты измерений}

% --- Таблица 1: Зависимость продольного напряжения от температуры ---
\subsection{Таблица 1. Зависимость продольного напряжения $U_{12}$ от температуры ($I = 1$ мА)}
\begin{center}
\begin{tabular}{|c|c|c|c|}
\hline
$T, K$ & 310 & 335 & 360 \\ \hline
$U_{12}, B$ & 0.10 & 0.26 & 0.31 \\ \hline
$1/T, 10^{-3} K^{-1}$ & 3.23 & 2.99 & 2.78 \\ \hline
$\ln \sigma$ & \dots & \dots & \dots \\ \hline
\end{tabular}
\end{center}

% --- Таблица 3: Зависимость ЭДС Холла от тока ---
\subsection{Таблица 3. Зависимость ЭДС Холла $U_x$ от тока ($B = 200$ мТл, $T = 300$ К)}
\begin{center}
\begin{tabular}{|c|c|c|c|c|}
\hline
$I, \mu A$ & 393 & 598 & 808 & 1008 \\ \hline
$U'_{34}, B$ & 41 & 62 & 84 & 105 \\ \hline
$U''_{34}, B$ & 42 & 64 & 83 & 107 \\ \hline
$U_x, B$ & \dots & \dots & \dots & \dots \\ \hline
\end{tabular}
\end{center}

\section{Пример расчета}
Для $I = 1008 \mu A$ и $B = 200$ мТл:
\begin{equation}
    U_x = \frac{105 - 107}{2} = -1.0 \text{ В (условно)}
\end{equation}
\textit{Примечание: В расчетах необходимо использовать толщину образца $d$ и геометрические параметры из паспорта установки.}

\section{Графики}
В данной работе строятся следующие зависимости:
\begin{enumerate}
    \item $\ln \sigma = f(1/T)$ --- для определения энергии активации.
    \item $U_x = f(I)$ --- для проверки линейности эффекта Холла.
    \item $U_x = f(T)$ --- для изучения температурной зависимости концентрации.
\end{enumerate}

\section{Выводы}
В ходе выполнения работы был исследован эффект Холла в полупроводнике. По знаку ЭДС Холла был определен тип проводимости (n-тип или p-тип). Рассчитаны основные характеристики носителей заряда: концентрация $n$ и подвижность $\mu$. Линейная зависимость $U_x(I)$ подтверждает теоретические положения.

\end{document}